
%% ============================================================
%% ВСТАВКА (23 октября 2026): масштаб G через вектор Бюргерса.
%% Честный статус: тождество-переформулировка, не вывод.
%% Вставить перед \section{Заключение}.
%% ============================================================

\section{Масштаб гравитации и вектор Бюргерса кольца}
\label{sec:burgers-scale}

\subsection{Тождество для гравитационной постоянной тонкой структуры}

В модели кольца конечного радиуса $a$ с вектором Бюргерса $b$
безразмерная сила гравитационно-подобного канала (Вариант~2)
управляется отношением $(b/a)^2$. Если отождествить радиус кольца с
комптоновской длиной электрона $a=\hbar/(m_ec)$, а вектор Бюргерса --- с
планковской длиной $b=\ell_{P}=\sqrt{\hbar G/c^3}$, то получается точное
алгебраическое тождество
\begin{equation}
\left(\frac{\ell_P}{a}\right)^2
= \frac{\hbar G/c^3}{\big(\hbar/m_ec\big)^2}
= \frac{G m_e^2}{\hbar c}
= \alpha_G \approx 1{,}75\times10^{-45},
\label{eq:burgers-identity}
\end{equation}
проверенное численно с относительной точностью $1{,}8\times10^{-16}$.
Требуемый вектор Бюргерса при этом равен $|b|=a\sqrt{\alpha_G}
=1{,}0000\,\ell_P$.

\subsection{Честный статус: переформулировка, а не вывод}

Необходимо явно подчеркнуть, во избежание ложного оптимизма: тождество
(\ref{eq:burgers-identity}) \textbf{не является независимым выводом}
величины $G$. Подстановка $b=\ell_P$ и $a=\hbar/(m_ec)$ использует $G$ и
$\hbar$ уже в определении $\ell_P$, поэтому равенство
$(\ell_P/a)^2=\alpha_G$ есть просто переписывание определения $\alpha_G$
в планковских единицах.

Проверка через энергетическое замыкание это подтверждает: попытка
зафиксировать $b$ из условия равенства собственной энергии
дислокационного кольца энергии покоя электрона,
$m_ec^2 = \gamma_{3D}\,b^2 a$ с $\gamma_{3D}=c^7/(8\pi\hbar G^2)$, даёт
$b\sim10^{-57}$~м, то есть $b/\ell_P\to0$, а комбинация
$\gamma_{3D}\ell_P^2 a / (m_ec^2)$ воспроизводит ровно фактор
$1/\alpha_G$ (расхождение в 44 порядка). Иными словами, попытка вывести
$b=\ell_P$ из планковской жёсткости \enquote{ходит по кругу}: она не
выделяет планковскую длину, а возвращает ту же величину $\alpha_G$.

\subsection{Что это меняет и что остаётся}

Тождество (\ref{eq:burgers-identity}) полезно: оно показывает, что
\emph{если} топология кольца фиксирует $b=\ell_P$ по какой-то
независимой причине (например, квантование потока/вектора Бюргерса,
вытекающее из самой калибровочной структуры Лазара и не использующее
$G$ заранее), то масштаб гравитации получается автоматически.
Но само по себе оно переводит нерешённый вопрос \enquote{почему
гравитация так слаба} в эквивалентный \enquote{почему $b=\ell_P$}, не
закрывая его. Таким образом, одновременное объяснение кулона и
гравитации в модели колец остаётся на уровне непротиворечивой
калибровки (два независимых модуля $a_1,f_1$), а не предсказания,
ровно как и констатировано в разделе~\ref{sec:variant3} и Заключении.
